
\subsubsection{Interpreting Sparse Coupling Detection Capability}

The core finding---methodology can detect coupling limited to 2 dominant dimension pairs rather than distributed across all 10 pairs---demonstrates capability to distinguish sparse from broad patterns in synthetic validation. If validated on real benchmarks, sparse coupling would suggest trustworthiness dimensions are architecturally independent by default, with coupling emerging only where specific mechanisms overlap.

This capability would help refute two alternative hypotheses. The broad independence hypothesis predicts zero coupling across all dimension pairs. Synthetic validation shows coupling is detectable when designed into data (6 pairs, phi 0.33-0.40, p < 1e-13). The broad coupling hypothesis predicts pervasive co-occurrence across many pairs from shared architectural bottlenecks. The h-m2 FAIL result demonstrates methodology can detect sparsity: 0 models exhibit $\geq 3$ significant pairs in synthetic validation.

If validated, sparse coupling would reveal that multi-dimensional trustworthiness failures are not universal---models do not simply fail everywhere when stressed. Instead, vulnerabilities would cluster in specific subsystems detectable by this framework. Truthfulness and robustness would couple where both depend on calibration; fairness and safety would couple where both are shaped by value alignment training. Remaining dimension pairs would remain largely independent.

\subsubsection{Difficulty as Suppressor in Synthetic Data}

The h-m1 finding---effect size retention 86-161\%, with 5 of 6 cases showing partial phi exceeding raw phi---demonstrates methodology capability to control for difficulty confounds in synthetic data. The suppressor effect interpretation (difficulty masks true coupling) may not replicate with real benchmarks.

Difficulty was constrained as suppressor variable by design in synthetic data: instances failing due to high difficulty may pass on dimension-specific vulnerabilities, and vice versa. Difficulty variance adds noise orthogonal to coupling signal; removing this noise unmasks designed coupling strength. This aligns with the independence constraint (|correlation(difficulty, dimension)| < 0.2).

Alternative explanations remain untested: (1) synthetic constraint artifact---real benchmark difficulty may show different relationships, (2) statistical artifact---partial correlation may inflate effect sizes in small samples with weak predictors, (3) measurement error---confidence-based difficulty proxy may not capture true difficulty. Real-world validation with unconstrained difficulty distributions is required.

\subsubsection{Limitations}

\textbf{Synthetic Data (High Impact).} All experiments used synthetic coupling data instead of real benchmarks. The methodology was validated---phi coefficient, partial correlation, Mantel test implementations work correctly---but coupling patterns in real LLMs remain uncharacterized. Synthetic data was designed with target coupling values (phi 0.35-0.40 for dominant pairs), so results confirm ''coupling is measurable if it exists'' rather than ''coupling exists in practice.''

MultiTrust and TrustLLM benchmarks remain access-gated. Real GPT-4, Claude-3, and Llama-3 API evaluations were not conducted. The contribution is methodological---demonstrating a pipeline for detecting and characterizing coupling. Real-world characterization requires production benchmark access.

\textbf{Sample Size (Medium Impact).} Mantel test used n=100 instances per dimension, insufficient for statistical power (requires n $\geq$ 500 per Legendre and Fortin, 2010). Criterion r < 0.7 was met (coupling matrices structurally dissimilar), but p-values remained non-significant (0.317-0.758).

Qualitative patterns in synthetic data are consistent---distinct coupling patterns observable across GPT-4, Claude-3, and Llama-3 profiles. Statistical inconclusiveness reflects sample size limitation, not absence of effect. This is standard proof-of-concept trade-off: demonstrate feasibility at small scale, power appropriately for confirmatory phase.

\textbf{Bonferroni Over-Correction (Low Impact).} h-m2 applied strict Bonferroni correction (alpha\_adj = 0.01/30 $\approx$ 0.00033), excluding borderline pairs. Sparse coupling conclusion is robust to correction method choice. Even if 1-2 additional pairs achieve significance under FDR control, no model would reach $\geq 3$ pairs threshold.

\subsubsection{Implications If Validated}

If validated on real benchmarks, coupling-aware model selection could reduce compound failures in safety-critical deployments. Deployment teams could prioritize models with strong coupling in required dimensions. Benchmark designers could ensure coverage of detected coupling pairs for compound risk assessment rather than treating dimensions independently.

The synthetic data limitation prevents premature deployment recommendations until real-world coupling patterns are validated. This work serves as methodology demonstration rather than actionable guidance.
